% ch5-values-effects.tex
\chapter{Values, Uses and Effects}
\label{ch:values}

\begin{answer}
An SSA (static single assignment) value in ulm.lam may be used any number
of times; using it twice means reading one matrix twice, not copying it.
Rewrites that would duplicate a \emph{computation} check that the value
has a single use. Whether an operation may be removed, moved out of a loop or
merged with another is decided by MLIR's generic passes from per-op effect
declarations: all \code{ulm_lam} ops are \code{Pure}, \code{ulm_lamk} ops
have no effects on tensors and read/write effects on memrefs, and a call of
an external function may do anything. Moving a pure op out of a loop is
sound for the declared facts. If a declaration is false, the move becomes
visible: a factorization hoisted out of a loop that runs zero times can
still report failure.
\end{answer}

% ---------------------------------------------------------------------------
\section{One value, several uses}
\label{sec:values-uses}

In \cref{fig:gls-dag}, \(W = L_1\inv X\) has two uses and \(L_1\) has two
uses. Each is computed once and stored once (\cref{sec:gls-storage}).
Nothing in MLIR restricts how often an SSA value is used, and nothing in
ulm.lam needs a restriction.

The restriction appears where a rewrite would copy an \emph{expression}
into several places. \code{Flatten} merges an inner product into the outer
one only if the inner product has no other use
(\srcref{lib/UlmLamDialect.cpp}{578}{602}):

\begin{lstlisting}[style=cpp]
for (Value v : op->getOperands()) {
    auto inner = v.template getDefiningOp<OpT>();
    if (inner && inner->hasOneUse()) {           // mul(a, mul(b, c)) -> mul(a, b, c)
        llvm::append_range(operands, inner->getOperands());
        changed = true;
    } else {
        operands.push_back(v);
    }
}
\end{lstlisting}

If \(P = AB\) is used in \(PC\) and in \(P + C\), flattening \(PC\) into
\(ABC\) would compute \(AB\) a second time for the sum. With the check,
\(P\) stays a value of its own. \code{examples/variants.jl} has exactly
this program (\code{shared}); \code{@code_lam} shows \(AB\) computed once:

\begin{textout}
shared(A::n×k, B::k×n, C::n×n)
  gemm     T1 ← A B
  gemm     T2 ← T1 C
  copy     T3 ← T1
  axpy     T3 ← T3 + C
  return T2, T3
\end{textout}

(The copy has another reason: the sum is accumulated in place into a new
matrix that is returned in Julia's column-major layout, so \(T_1\) is
copied there transposed before \(C\) is added. It is not a copy for the
second use; \(T_1\) is read, not overwritten, by the first.)
\code{PushTranspose}, which turns \((ABC)\T\) into \(C\T B\T A\T\), has the
same one-use condition (\srcref{lib/UlmLamDialect.cpp}{641}{677}).

\begin{theory}[ (\FGs{1.1}, \FGs{2.7})]
In the terms of \FG{1.1}{Structural rules and substructural type systems},
SSA values obey both \emph{contraction} (a value may be used twice) and
\emph{weakening} (a value may be unused): MLIR's value discipline is the
ordinary, non-substructural one. That is safe because tensors have value
semantics, \FG{2.7}{Value semantics vs reference/memory semantics in IRs}:
a use can never observe another use. The one-use checks above are not a
typing rule; they are a cost side condition on rewrites, so that a rewrite
never duplicates work.
\end{theory}

\begin{limit}
An SSA value with several uses does not imply several copies of its
storage. Copies are made by bufferization, and only when a use would
\emph{overwrite} storage that another use still reads (\cref{ch:storage}).
Conversely, linear types (\FGs{1.2}) would restrict the number of uses;
nothing in MLIR or ulm.lam does.
\end{limit}

% ---------------------------------------------------------------------------
\section{Effects decide what may move}
\label{sec:effects}

\begin{steps}
\item[Input] an op inside a loop, or an op whose results are unused.
\item[Representation] the op's \code{MemoryEffectOpInterface}: the trait \code{Pure} (no effects, speculatable), or a \code{getEffects} method.
\item[Rule] generic passes: DCE (dead code elimination) removes ops without effects whose results are unused; LICM (loop-invariant code motion) hoists pure ops whose operands are defined outside the loop; CSE merges identical pure ops.
\item[Output] the program with ops removed, moved or merged.
\item[Theory] \FG{1.6}{Effect systems and algebraic effects}; \FG{3.2}{Traits and interfaces}.
\end{steps}

\paragraph{The declarations}
\begin{itemize}
\item Every \code{ulm_lam} op is \code{Pure} (the base class
  \code{UlmLam_Op}, \srcref{include/ulm_lam/UlmLamOps.td}{23}{35}). They
  are mathematical functions of their operands.
\item \code{ulm_lamk} ops implement \code{getEffects} with
  \code{dpsEffects} (\srcref{lib/UlmLamDialect.cpp}{808}{824}): memref
  operands are read, DPS destinations are also written; ``Tensor operands
  have value semantics (no effects).'' The same op is effect-free on tensors
  and effectful on memrefs.
\item \code{ulm_lamk.call} calls a function written in Julia or C; it adds
  a read and a write of the default resource, i.e.\ of any memory
  (\srcref{lib/UlmLamDialect.cpp}{1150}{1160}): ``the call is never
  hoisted, duplicated or removed''.
\end{itemize}

\paragraph{Removing} On tensors, a kernel whose results are unused is dead,
even a factorization with its \code{info}: in \code{v12-dead-potrf.mlir}
the canonicalizer erases the tensor \code{potrf} and keeps the memref one
(\code{outputs/dead.out}):

\begin{textout}
func.func @f(%arg0: tensor<4x4xf64>, %arg1: memref<4x4xf64>) -> tensor<4x4xf64> {
  %info = ulm_lamk.potrf uplo(lower) outs(%arg1 : memref<4x4xf64>)
  return %arg0 : tensor<4x4xf64>
}
\end{textout}

A run-time check that is a by-product of a computation lives only as long
as the computation is needed. On tensors this is correct by construction:
\code{info} is an ordinary result, and an unused result is not observable.

% ---------------------------------------------------------------------------
\section{Moving a factorization out of a loop}
\label{sec:effects-licm}

The loop in \code{examples/variants.jl} applies \(M\inv\) in every step:

\begin{lstlisting}[style=jl]
@lam function iterate(M::Mat(:n, :n, :spd), x0::Vec(:n), steps::Index)
    loop(steps, x0) do i, x
        (inv(M) * x,)
    end
end
\end{lstlisting}

The traced IR has the inverse inside the loop body; after normalization,
inference and \code{loop-invariant-code-motion} it is outside
(\code{outputs/iterate-traced.mlir}, \code{outputs/iterate-licm.mlir}):

\begin{lstlisting}
// traced
%v4 = scf.for %v5 = %v2 to %p.steps step %v3 iter_args(%v6 = %p.x0) -> (tensor<?x1xf64>) {
  %v7 = ulm_lam.inverse %v1 : tensor<?x?xf64> -> tensor<?x?xf64>
  %v8 = ulm_lam.mul %v7, %v6 : tensor<?x?xf64>, tensor<?x1xf64> -> tensor<?x1xf64>
  scf.yield %v8 : tensor<?x1xf64>
}
// after LICM
%0 = ulm_lam.inverse %arg1 : ... -> tensor<?x?xf64, #ulm_lam.props<square, symmetric, spd, ...>>
%1 = scf.for %arg4 = %c0 to %arg3 step %c1 iter_args(%arg5 = %arg2) -> (tensor<?x1xf64>) {
  %2 = ulm_lam.mul %0, %arg5 : ...
  scf.yield %2 : tensor<?x1xf64>
}
\end{lstlisting}

LICM moved \code{inverse} because its operand is defined outside the loop
and the op is \code{Pure}. Kernel selection then factorizes \(M\) once,
before the loop; the loop body is two \code{trsv}
(\code{outputs/iterate-selected.mlir}). Kernel selection would place the
factorization there even without LICM, because it factorizes a matrix
directly after its definition (\cref{sec:gls-select}); two independent
mechanisms agree.

% ---------------------------------------------------------------------------
\section{Speculation and a false declaration}
\label{sec:effects-speculation}

\code{Pure} in MLIR means \code{NoMemoryEffect} \emph{and}
\code{AlwaysSpeculatable}: the op may be executed even on paths where the
original program would not execute it. For \code{ulm_lam.inverse} of an SPD
matrix this is harmless: the inverse exists, and computing its
factorization early changes nothing observable. For an \(M\) that is
declared SPD but is not, it changes the outcome
(\code{examples/speculate.jl}, \code{outputs/speculate.out}):

\begin{textout}
steps = 0: InfoError: iterate: info 1: Cholesky factorisation of v2: not positive definite
steps = 1: InfoError: iterate: info 1: Cholesky factorisation of v2: not positive definite
\end{textout}

With \code{steps = 0} the formula never applies \(M\inv\); its value is
\(x_0\). The compiled function factorizes \(M\) before the loop and reports
the failure anyway. Relative to the declared facts (\(M\) is SPD) the
transformation is exact; relative to the actual input it is not.

\begin{theory}[ (\FGs{1.6}, \FGs{3.2})]
\FG{1.6}{Effect systems and algebraic effects} explains why compilers need
effect information~\cite{effects}: a computation without effects may be removed, hoisted
or merged. MLIR attaches that information to op classes and instances
through traits and interfaces, \FG{3.2}{Traits and interfaces}, and LICM,
CSE and DCE are written once against them. The example shows the other
half of the contract: correctness depends on the declarations being honest
\emph{for the inputs that occur}. \code{ulm_lam.inverse} is pure on the
domain where the inverse exists; the domain is guaranteed only by a
declaration.
\end{theory}

\begin{limit}
Failure is not modeled as an effect. In an effect system,
``may fail'' would be an effect like ``may raise'', and a may-fail
operation would not be speculatable. In ulm.lam the possibility of failure
is outside the \code{ulm_lam} level entirely (the ops are total over the
declared domain); it appears at the \code{ulm_lamk} level as an ordinary
result value, \code{info}, and only \code{ulm-lam-info} turns it into
control flow (\cref{sec:effects-info}).
\end{limit}

% ---------------------------------------------------------------------------
\section{Failure as a value, then as control flow}
\label{sec:effects-info}

\code{ulm_lamk.potrf} returns \code{info} as an \code{i64} result. After
bufferization, \code{ulm-lam-info} (\srcref{lib/UlmLamInfo.cpp}{1}{41})
numbers every check of a function and rewrites the rest of the block into
an \code{scf.if}:

\begin{lstlisting}
%info = ulm_lamk.potrf ... outs(%A)
%code = arith.select (%info != 0), k, 0
%s = scf.if (%code == 0) -> i64 { <rest>; yield <status of rest> }
                         else { yield %code }
\end{lstlisting}

(This listing is the schematic from the pass's header comment.) A loop
carries the status, so the iterations after a failure do nothing. The
function returns the status; outputs after a failure are unspecified, as in
LAPACK. The descriptions of the checks go into the function attribute
\code{ulm_lam.info} and from there into the C++ interface comment and
Julia's \code{InfoError}.

\begin{theory}
This is error handling by explicit status passing: every operation that can
fail returns a code, and the compiler threads it through the rest of the
computation. An exception or an algebraic effect with a handler
(\FGs{1.6}) would express the same thing with a non-local exit;
\code{scf.if} nesting is how a structured IR without exceptions encodes it.
\end{theory}

% ---------------------------------------------------------------------------
\section{Regions in MLIR and regions in memory management}
\label{sec:values-regions}

ulm.lam's IR contains MLIR \emph{regions}: the body of \code{scf.for}, the
branches of \code{scf.if}, the body of every \code{func.func}. These are
nested blocks of code with lexical scoping of SSA values,
\FG{2.6}{Regions and structured control flow in the IR}.

\begin{limit}
MLIR regions are not the regions of region-based memory management,
\FG{1.5}{Region types and lifetimes}, where a region is a pool of memory
whose lifetime the type system tracks. The closest thing in ulm.lam is the
workspace (\cref{sec:storage-workspace}): one block of memory per call, from
which all temporaries are carved and which is released (or kept by the
caller) at the end. That is an arena with a lifetime of one call, managed
by a pass and a calling convention, not by types.
\end{limit}
